
\documentclass[a4paper,12pt]{article}

\usepackage{JYM-harj}
\usepackage{tikz}
\usepackage{amsmath}

\setlength{\parindent}{0in}

\setlength{\voffset}{-1.0cm}
\addtolength{\textheight}{2.0cm}

\setlength{\hoffset}{-1.0cm}
\addtolength{\textwidth}{2.0cm}

\usepackage{graphicx}
\usepackage{verbatim}

\usepackage{hyperref}

%\usepackage[finnish]{babel}
%\usepackage[utf8]{inputenc}
%\usepackage[T1]{fontenc}
%\usepackage{amsmath,amsfonts,amssymb,amsthm,amscd}


% 1: Teht\"av\"at
% 2: Kaikki ratkaisut
% 3: T\"ahdett\"om\"at ratkaisut
% 4: T\"ahdelliset ratkaisut
\newcommand{\tversio}{1}

\newcommand{\harjoitusnro}{3.3.2026 klo 17:00}

\ifthenelse{\tversio = 1}{\pagestyle{empty}}{}

\begin{document} 

\otsikko

\begin{enumerate}

\item 
    \begin{enumerate}
    \end{enumerate}
    
\item 
\begin{enumerate}
\end{enumerate}

\item Todista induktiolla, että $\frac{1}{1\cdot 2}+\frac{1}{2\cdot 3}+\frac{1}{3\cdot 4}+\cdots+\frac{1}{n(n+1)}=\frac{n}{n+1}$ kaikilla $n \geq 1$. (6p) \\

\textbf{Vastaus:}

\underline{Alkuaskel:} Pienin mahdollinen arvo muuttujalle $n$ on $1$. Katsotaan pitääkö väite paikkansa jos $n=1$:
\begin{equation*}
    \frac{1}{1\cdot(1+1)}=\frac{1}{1\cdot 2}=\frac{1}{2}
\end{equation*}
ja
\begin{equation*}
    \frac{1}{1+1}=\frac{1}{2}
\end{equation*}
Yhtälön vasen ja oikea puoli ovat yhtä suuret, eli väite pätee kun $n=1$. 

\underline{Induktio-oletus:} Oletetaan, että väite pätee mielivaltaisella kokonaisluvulla $k\geq1$. Oletetaan, että seuraava yhtälö pätee:
\begin{equation*}
    \frac{1}{1\cdot 2}+\frac{1}{2\cdot 3}+\frac{1}{3\cdot 4}+\cdots+\frac{1}{k(k+1)}=\frac{k}{k+1}
\end{equation*}

\underline{Induktioväite:} Osoitetaan, että jos oletus pätee muuttujalla $k$, niin se pätee myös seuraajalla $k+1$, ja että väite pätee myös muodossa:

\begin{equation*}
       \frac{1}{1\cdot 2}+\frac{1}{2\cdot 3}+\frac{1}{3\cdot 4}+\cdots+\frac{1}{k(k+1)}+\frac{1}{(k+1)((k+1)+1)}=\frac{k+1}{(k+1)+1}
\end{equation*}

Yllä olevan yhtälön oikea puoli sievenee  muotoon $\frac{k+1}{(k+1)+1}=\frac{k+1}{k+2}$.

\underline{Induktioväitteen todistus:} Induktioväitteen yhtälön vasemman puolen alku on tismalleen sama kuin induktio-oletuksen yhtälön vasen puoli: $\frac{1}{1\cdot 2}+\frac{1}{2\cdot 3}+\frac{1}{3\cdot 4}+\cdots+\frac{1}{k(k+1)}$. Lähdetään muokkaamaan induktioväitteen vasenta puolta:
\begin{align*}
    \frac{k}{k+1}+\frac{1}{(k+1)(k+2)}&=\frac{k(k+2)}{(k+1)(k+2)}+\frac{1}{(k+1)(k+2)}\\
    &=\frac{k^2+2k+1}{(k+1)(k+2)} \\
    &=\frac{(k+1)^2}{(k+1)(k+2)} \\
    &=\frac{k+1}{k+2} \\
\end{align*}
Lopputulos vastaa induktioväitteen oikeaa puolta, joten väite pätee.

\underline{Johtopäätös:} Alkuaskeleesta ja induktioaskeleesta seuraa induktioperiaatteen nojalla, että $\frac{1}{1\cdot 2}+\frac{1}{2\cdot 3}+\frac{1}{3\cdot 4}+\cdots+\frac{1}{n(n+1)}=\frac{n}{n+1}$ kaikilla $n \geq 1$.

\end{enumerate}

\end{document}